\documentclass[11pt]{article}
\usepackage[a4paper,margin=23mm]{geometry}
\usepackage[T1]{fontenc}
\usepackage{lmodern,amsmath,amssymb,amsthm,mathtools,microtype,booktabs}
\usepackage[colorlinks=true,linkcolor=blue,citecolor=blue,urlcolor=blue]{hyperref}
\newcommand{\E}{\mathbb E}
\newcommand{\Tr}{\operatorname{Tr}}
\newcommand{\tr}{\operatorname{tr}}
\newcommand{\HH}{\mathcal H}
\newcommand{\WW}{\mathcal W}
\newcommand{\dd}{\mathrm d}
\newcommand{\norm}[1]{\left\lVert#1\right\rVert}
\newtheorem{theorem}{Theorem}
\newtheorem{lemma}{Lemma}
\title{Polynomial Slater preparations and uniform orbit densities\\
\large Exact neutral-slope pinching on the selected polynomial core}
\author{Alan Day \and Chengheng\thanks{AI research assistant (OpenAI).}}
\date{8 November 2026\quad Local research draft}
\begin{document}\maketitle
\begin{abstract}
We extend the finite-$N$ linear preparation estimate to fixed Gaussian
degree in the half-filled double-scaled regime. Partial traces of the
coefficient seed have weights within a controlled distance of a Slater
weight. A Weyl-orbit count bounds the dimensions of every supporting
irreducible representation. It controls all oriented tensor flattenings
of the orbit synthesis map. Iterated fixed-Schatten Gaussian inequalities
and a conditional Wick identity then give dimension-uniform density
moments. This transfers exact energy-pinching errors to every fixed
selected polynomial preparation. Two quadratic preparations give explicit
normalized slope correlators. Multiple interlaced nonpolynomial pinched
factors are not inferred from convergence on this core.
\end{abstract}

\section{Setting and quantitative statement}
Use standard CAR and the same overlapping Hermitian complex Gaussian
$p$-body couplings as in the preceding Slater drafts. Set
$m=N/2$, $d_m=\binom Nm$, $\HH_m=\bigwedge^m\mathbb C^N$ and fix a
Slater determinant $O$. Real independent standard Gaussian coordinates
$g_a$ carry the orthogonal noise representation $R(U)$ of $U(N)$.
Let $P_O(g)$ have degree at most $K\ge1$ and obey
\begin{equation}\label{stab}
 P_O(k\cdot g)=k_mP_O(g)k_m^\dagger,
 \qquad k\in U(m)\times U(N-m).
\end{equation}
Write $P_OO=\sum_{k=0}^K P_{k,O}O$ for its orthogonal Wick-chaos
decomposition and $a_k=\E\norm{P_{k,O}O}^2$, $a=\sum_k a_k$.
All these quantities may depend on $N$.

\begin{theorem}\label{main}
Assume $p^2/N\le\Lambda$ and $p\le N/(12K)$, and put
\begin{equation}\label{constants}
 C_K=e^{36K^2\Lambda+1},\qquad B_{0,r}=1,\qquad
 B_{k,r}=\frac{[k(2r-1)]^k}{k!}C_K^{1-1/r}\quad(k\ge1).
\end{equation}
There is a positive covariant orbit density $R_P$ such that for covariant
$Y$ in joint Schatten $L^2$,
\begin{align}
 \E\langle O|P_O^\dagger YP_O|O\rangle&=\E\tr_m(R_PY),\label{identity}\\
 \norm{R_P}_{L^r(\tr_m)}&\le
 \left(\sum_{k=0}^K\sqrt{B_{k,r}a_k}\right)^2
 \le (K+1)\max_{k\le K}B_{k,r}\,a .\label{bound}
\end{align}
Here $r\ge1$ is an integer and $\tr_m=\Tr/d_m$.
\end{theorem}
The constants are not optimized. For fixed $K$ and $p^2/N\to\lambda>0$,
the hypotheses hold eventually. This is a fixed-degree statement,
not a bound uniform in degree or a general matrix-polynomial
hypercontractivity assertion.

\newpage
\section{Orbit synthesis and the coefficient seed}
Normalized Haar measure is understood. Define
\begin{equation}\label{synthesis}
 q_U(g)=U_m^\dagger P_O(U\cdot g)O,\qquad
 Z(g)f=\sqrt{d_m}\int q_U(g)f(U)\,\dd U,\qquad R_P=ZZ^\dagger.
\end{equation}
Only a finite-dimensional space of matrix coefficients is needed in the
Haar domain. Gaussian invariance and $Y(U\cdot g)=U_mY(g)U_m^\dagger$
give \eqref{identity}. They also give $\E\tr_m R_P=a$.
Chaos projection commutes with orthogonal noise rotations, so
\eqref{stab} holds separately in each chaos.

For degree $k$, use a symmetric coefficient tensor in ordered indices:
\begin{equation}\label{seed}
 P_{k,O}(g)O=\sum_{i_1,\ldots,i_k} A_{i_1\cdots i_k}
 :g_{i_1}\cdots g_{i_k}:,\qquad
 \Psi=\sum_{\boldsymbol i}e_{i_1}\otimes\cdots\otimes e_{i_k}
                     \otimes A_{\boldsymbol i}.
\end{equation}
Wick isometry in this ordered symmetric convention reads
\begin{equation}\label{isometry}
 a_k=k!\norm{\Psi}^2.
\end{equation}
Let $\rho=R^{\otimes k}\otimes\bigwedge^m$. The coefficient tensor
of $q_U$ is $\Psi_U=\rho(U^{-1})\Psi$. For a diagonal unitary $D$,
stabilizer covariance and $D_mO=\chi_O(D)O$ imply
\begin{equation}\label{weight}
 \rho(D)\Psi=\chi_O(D)\Psi,\qquad
 w=(\underbrace{1,\ldots,1}_{m},0,\ldots,0).
\end{equation}
This is a statement about the full tensor, before taking any partial
trace.

Complexify the real Hermitian noise space and choose an orthonormal
weight basis. It is unitarily equivalent to the conjugate adjoint
representation on $\operatorname{End}(\bigwedge^p\mathbb C^N)$.
Its weights are differences of two $p$-subset indicators. Every noise
weight $\nu$ has integral entries, sum zero and $\norm{\nu}_1\le2p$.
This change of coefficient basis does not change any flattening's
singular values; we still use real independent Gaussians for the
probabilistic estimates.

Trace out a set $J$ of $\ell$ noise legs. In a product weight basis,
\begin{equation}\label{partial}
 T_I=\Tr_{\WW^{\otimes J}}|\Psi\rangle\langle\Psi|
     =\sum_\beta|\psi_\beta\rangle\langle\psi_\beta|,
 \quad I=[k]\setminus J.
\end{equation}
Each nonzero $\psi_\beta$ has a weight $z=w-\nu_\beta$ with
\begin{equation}\label{near}
 \sum_i z_i=m,\qquad\norm{z-w}_1\le2\ell p.
\end{equation}
Thus partial tracing loses the exact Slater weight but retains a
quantitative neighborhood of it. This is the additional input needed
beyond the preceding linear argument.

\newpage
\section{A near-Slater Weyl-orbit bound}
\begin{lemma}\label{weyl}
Under the hypotheses of Theorem \ref{main}, every integral weight
satisfying \eqref{near}, with $0\le\ell\le K$, has at least
$d_m/C_K$ distinct coordinate permutations. Every irreducible
$U(N)$ representation containing such a weight therefore has
dimension at least $d_m/C_K$.
\end{lemma}
\begin{proof}
For $\ell=0$ the orbit has exactly $d_m$ elements. Otherwise let $b$
be the number of entries of $z$ outside $\{0,1\}$ and let $S$ be
their sum. Since $z-w$ is integral of total sum zero, its total
positive and negative masses are both at most $\ell p$. Consequently
\begin{equation}
 b\le2\ell p,\quad -\ell p\le S\le b+\ell p,\quad
 |S|,|S-b|\le3\ell p=:D.
\end{equation}
The counts of ones and zeros are $n_1=m-S$, $n_0=m-b+S$.
Put $\delta_1=-S$, $\delta_0=S-b$; then
$|\delta_i|\le D\le m/2$ and $\delta_0+\delta_1=-b$.

For an integer $|\delta|\le m/2$,
\begin{equation}\label{factorial}
 \log\frac{(m+\delta)!}{m!}
 \le\delta\log m+\frac{\delta^2+|\delta|}{m}.
\end{equation}
For positive $\delta$, expand the logarithm as
$\delta\log m+\sum_{j=1}^{\delta}\log(1+j/m)$ and use
$\log(1+x)\le x$. For $\delta=-t$, expand it as
$-t\log m-\sum_{j=0}^{t-1}\log(1-j/m)$ and use
$-\log(1-x)\le2x$ for $0\le x\le1/2$. Both give
\eqref{factorial}, including the endpoints.

If $n_t$ counts the entries equal to $t\notin\{0,1\}$, the orbit
cardinality $L(z)$ satisfies
\begin{align}
 \frac{L(z)}{d_m}
 &=\frac{m!^2}{n_0!n_1!\prod_{t\ne0,1}n_t!}\\
 &\ge\frac{m^b}{b!}
       \exp\left[-\frac{2D^2+2D}{m}\right]
 \ge\exp[-36\ell^2\Lambda-1]\ge C_K^{-1}.\label{orbitbound}
\end{align}
We used $\prod n_t!\le b!$, $b\le m$, and
$2D/m=12\ell p/N\le1$.
Permutation matrices implement the Weyl group. Their distinct weights
have orthogonal weight spaces in a unitary representation, so its
dimension bounds their number from below.
\end{proof}

No multiplicity-free assumption is involved. The estimate permits
entries such as $2$ and $-1$ as well as imbalanced numbers of zeroes
and ones. It is deliberately coarse; the bounded $p^2/N$ and fixed
$K$ assumptions, rather than closeness in operator norm, keep its
loss independent of $N$.

\newpage
\section{All oriented flattenings and decoupled moments}
Let $F_I$ flatten the synthesis coefficient tensor with physical
output and noise legs $I$ on its row side, and the Haar input and
remaining noise legs $J$ on its column side. Summing over the column
noise indices cancels their unitary rotations. Thus
\begin{equation}\label{frame}
 F_IF_I^\dagger=d_m\int\rho_I(U^{-1})T_I
                    \rho_I(U^{-1})^\dagger\,\dd U,
 \qquad \norm{F_I}_{\mathrm{HS}}^2=d_m\norm{\Psi}^2.
\end{equation}
Decompose the row representation as $\bigoplus_\alpha V_\alpha
\otimes\mathcal M_\alpha$. Every supporting irreducible component
of $T_I$ contains one of the weights in \eqref{near}: project each
$\psi_\beta$ in \eqref{partial} onto the isotypic components, using
that these projections commute with the torus. Schur averaging gives
\begin{equation}\label{schur}
 F_IF_I^\dagger=\bigoplus_\alpha\frac{d_m}{\dim V_\alpha}
 I_{V_\alpha}\otimes\Tr_{V_\alpha}(T_{I,\alpha\alpha})
 \le C_K\norm{\Psi}^2 I.
\end{equation}
Here each multiplicity matrix is positive with norm no larger than
its trace, and $\Tr T_I=\norm{\Psi}^2$. Hence, for all orientations,
\begin{equation}\label{flattenmom}
 \norm{F_I}_{S_{2r}}^{2r}
 \le d_m C_K^{r-1}\norm{\Psi}^{2r}.
\end{equation}

For clarity we derive the exact probabilistic form required here.
For $X=\sum_i g_i A_i$, Hermitian dilation and the fixed-Schatten
Gaussian inequality give
\begin{equation}\label{rect}
 \E\norm{X}_{S_{2r}}^{2r}\le\frac{(2r-1)^r}{2}
 \left[\Tr\left(\sum_i A_iA_i^\dagger\right)^r
       +\Tr\left(\sum_i A_i^\dagger A_i\right)^r\right].
\end{equation}
Indeed, Gaussian integration by parts bounds $\E\Tr S^{2r}$ by
$(2r-1)\E\Tr(V|S|^{2r-2})$. Trace H\"older followed by scalar
H\"older gives $(\E\Tr S^{2r})^{1/r}\le
(2r-1)(\Tr V^r)^{1/r}$; dilation contributes the factor $1/2$.
This is the standard matrix-Gaussian argument \cite{tropp}.

Let $Z_{\rm dec}$ use independent Gaussian families for its $k$ slots.
Condition on $k-1$ families and apply \eqref{rect}. The two terms
are the $2r$-moments of a row concatenation and a column stacking.
Repeat in each branch. Each noise slot is assigned once to a row or
column, producing precisely all $2^k$ oriented flattenings:
\begin{align}
 \E\norm{Z_{\rm dec}}_{S_{2r}}^{2r}
 &\le(2r-1)^{kr}2^{-k}\sum_{I\subset[k]}
                      \norm{F_I}_{S_{2r}}^{2r}\label{iteration}\\
 &\le d_m(2r-1)^{kr}C_K^{r-1}(a_k/k!)^r .\nonumber
\end{align}
The flattening iteration is established methodology \cite{chaos};
we retain fixed Schatten moments, not a logarithmic-dimensional
operator-norm specialization.

\newpage
\section{Coupled Wick chaos and sums of degrees}
Let $g^{(1)},\ldots,g^{(k)}$ be independent standard Gaussian vectors
and let $g=k^{-1/2}\sum_jg^{(j)}$. Conditional means are $g/\sqrt k$;
conditional cross-covariances between distinct slots are $-I/k$.
The conditional exponential generating function is
\begin{equation}\label{generating}
 \E\left[\exp\left(\sum_j t_j\cdot g^{(j)}\right)\middle|g\right]
 =\exp\left[\frac{g\cdot\sum_jt_j}{\sqrt k}
       +\frac12\sum_j|t_j|^2-\frac1{2k}\left|\sum_jt_j\right|^2\right].
\end{equation}
Differentiate once in each slot. Terms involving two derivatives of
the same $t_j$ cannot contribute. Comparing with the Wick generating
function yields, including repeated coordinate indices,
\begin{equation}\label{conditional}
 \E\left[\prod_{j=1}^k g^{(j)}_{i_j}\middle|g\right]
 =k^{-k/2}:g_{i_1}\cdots g_{i_k}:.
\end{equation}
Consequently $Z_k(g)=k^{k/2}\E[Z_{\rm dec}\mid g]$ exactly. Convexity
of the $2r$-th power of a Schatten norm gives
\begin{equation}
 \E\norm{Z_k}_{S_{2r}}^{2r}
 \le k^{kr}\E\norm{Z_{\rm dec}}_{S_{2r}}^{2r}.
\end{equation}
With \eqref{isometry} and \eqref{iteration}, this proves
\begin{equation}\label{purebound}
 \left(d_m^{-1}\E\Tr (Z_kZ_k^\dagger)^r\right)^{1/r}
 \le B_{k,r}a_k.
\end{equation}
This proof handles Wick polynomials with repeated indices directly;
no square-free-chaos restriction is imported from an operator-norm
decoupling theorem.

The zero-chaos vector has weight $w$ in $\bigwedge^m\mathbb C^N$,
whose weight space is one-dimensional, so $P_{0,O}O=cO$. Its orbit
density is $|c|^2I$ and $B_{0,r}=1$ is exact. Since $Z=\sum_kZ_k$,
Minkowski in the Gaussian--Schatten space, normalized by $d_m$, gives
\begin{equation}
 \norm{R_P}_{L^r}^{1/2}
 =d_m^{-1/(2r)}\norm{Z}_{L^{2r}(S_{2r})}
 \le\sum_{k=0}^K\norm{R_{P_k}}_{L^r}^{1/2}.
\end{equation}
Cauchy--Schwarz and Gaussian-chaos orthogonality finish
\eqref{bound}. This argument retains interference between degrees;
it does not assert $R_P=\sum_k R_{P_k}$ samplewise.

For $k=1$, the preceding linear-stage estimate
$\E\tr_m R_{P_1}^r\le(2r-1)!!a_1^r$ is stronger and holds without
half filling or double scaling. The present constants are only used
to secure finite, $N$-independent bounds at every fixed degree.

\newpage
\section{Exact pinching on the polynomial core}
For a covariant $X_N$, \eqref{identity} and H\"older imply
\begin{equation}\label{transfer}
 \E\norm{X_NP_NO}^2
 \le\norm{R_{P,N}}_{L^2(\tr_m)}\norm{X_N}_{L^4(\tr_m)}^2.
\end{equation}
Take $P_N$ to be the evaluation of a fixed polynomial in the previously
constructed $h_N,u_N,M_N,M_N^\dagger$. Its Gaussian
degree is fixed, it obeys stabilizer covariance, and its total
$a_N$ stays bounded by the established same-$J$ fixed-word limits.
Thus any covariant error going to zero in tracial $L^4$ goes to zero
as a vector in every such preparation.

In particular, use the actual normalized unpinched slope from the
neutral-slope stage,
\begin{equation}
 t_N=\frac{m\epsilon\alpha(\Phi_m(H_N)-\bar\kappa_mH_N)}
 {8\lambda_N C_m^{3/2}q_N\sqrt{\lambda_N}},\qquad
 b_N=(1-\mathcal P_{H_N})t_N,
\end{equation}
where $h_N=H_N/\sqrt{C_m}$, $\lambda_N=p^2/N$,
$q_N=e^{-4\lambda_N}$, $\alpha=1-p/m$, $\epsilon=1-\alpha^2$.
The symbols $\Phi_m,\bar\kappa_m$ retain their exact covariance-channel
and regression definitions from that stage; $\mathcal P_{H_N}$ keeps
entire degenerate energy blocks. Its tracial results give
\begin{equation}
 \norm{\mathcal P_{H_N}t_N}_{L^r(\tr_m)}\longrightarrow0,
 \qquad \norm{t_N-u_N}_{L^r(\tr_m)}\longrightarrow0
 \quad\hbox{for every fixed finite }r.
\end{equation}
The second statement also follows by interpolating the rank/size $L^2$
bridge with the uniform higher moments. Equation \eqref{transfer}
therefore yields
\begin{equation}\label{core}
 \E\norm{(b_N-u_N)P_NO}^2\longrightarrow0.
\end{equation}
These are dimensionless string-energy observables with $\hbar=1$,
fixed $\lambda>0$ and annealed half-filled Slater expectations. No
thermal or gravitational interpretation is added by this estimate.

More generally, let $W_N$ be a fixed covariant block of $h_N,u_N,b_N$
and fixed-time $h_N$ phases. Tracial H\"older and telescoping show
$W_N(b_N)-W_N(u_N)\to0$ in $L^2$; fixed higher norms bound all other
factors. Apply \eqref{identity} and polarize $P_N,Q_N$ to obtain
\begin{equation}\label{block}
 \varphi_N\bigl(P_N^\dagger W_N(b_N)Q_N\bigr)
 \longrightarrow\varphi\bigl(P^\dagger W(h,u)Q\bigr).
\end{equation}
The unpinched limit on the right is the established Slater transfer
representation, not the tracial vacuum unless its preparation says so.
An arbitrary number of selected factors is allowed in the fixed
polynomials $P,Q$; the exact-pinched factors remain inside a covariant
block. This does not yet justify raw words with several such blocks
separated by selected operators: their intervening preparations are
no longer Gaussian polynomials. Nor does core convergence alone
establish those high moments.

\newpage
\section{Quadratic preparations and explicit slope responses}
In the subset labels below, $O$ and $E$ denote the occupied and
unoccupied one-particle index sets. Then
\begin{equation}
 M_N^\dagger=L_N^{-1/2}
 \sum_{A\subset E,\,B\subset O}J_{AB}C_A^\dagger C_B,
 \quad L_N=\binom mp\binom{N-m}p,
\end{equation}
with $|A|=|B|=p$ and GUE covariance
$\E J_{AB}J_{CD}=\delta_{AD}\delta_{BC}$. The normalization of $h_N$
removes the common physical coupling scale.
For $P_N=M_N^\dagger h_N$, a possible internal contraction pairs
the left selected label $(A,B)$ with the right label $(B,A)$.
The latter annihilates particles in $A\subset E$ and hence kills
$O$. For $P_N=(M_N^\dagger)^2$, a selected label's reverse is never
selected. Both preparation vectors consequently have zero chaos-zero
part and are pure second Wick chaos. This argument is finite $N$.

The established transfer rule, with
$\gamma=q^{1/4}$, $q=e^{-4\lambda}$ and $s=\sqrt q$, gives
\begin{align}
 a_{M^\dagger h}&=\varphi(hMM^\dagger h)=q(1+s),\\
 a_{(M^\dagger)^2}&=\varphi(M^2(M^\dagger)^2)=2q.\label{norms}
\end{align}
Both are positive at fixed $\lambda>0$. At small finite $N$, the
second vector can vanish when $2p>m$; no normalization is taken there.
Let the normalized preparation expectation divide by its finite-$N$
$a_N$. Applying \eqref{block} to $b_N^2$ and evaluating the existing
same-$J$ word rule yields
\begin{align}
 \lim_N\frac{\varphi_N(h_NM_Nb_N^2M_N^\dagger h_N)}
                  {\varphi_N(h_NM_NM_N^\dagger h_N)}
 &=1+\frac{\lambda s q(1+q)}{1+s},\label{quad1}\\
 \lim_N\frac{\varphi_N(M_N^2b_N^2(M_N^\dagger)^2)}
                  {\varphi_N(M_N^2(M_N^\dagger)^2)}
 &=1+2\lambda s q(1+q).\label{quad2}
\end{align}
For example, before normalization the first numerator is
$q[1+s+\lambda\gamma^6(1+q)]$, and the second is
$2q[1+2\lambda\gamma^6(1+q)]$. The code checks these word
polynomials in algebraically independent $\gamma$ and
$\sqrt\lambda$ before imposing the physical curve.

The exact finite fixtures retain the same disorder and CAR signs:
\begin{center}
\begin{tabular}{ccrr}
\toprule
$N$ & $p$ & $\E\norm{M_N^\dagger h_NO}^2$
          & $\E\norm{(M_N^\dagger)^2O}^2$\\\midrule
4&1&$2/3$&$1/2$\\
6&1&$11/12$&$8/9$\\
4&2&$1/6$&$0$\\
6&2&$1/6$&$0$\\
8&2&$2/9$&$1/18$\\\bottomrule
\end{tabular}
\end{center}
Here $m=N/2$. These are finite algebra checks, not fits to the
double-scaled asymptotic values in \eqref{norms}.

\newpage
\section{Evidence, attribution and continuation}
The exact driver checks the original finite CAR preparations,
symmetrizes their ordered Gaussian coefficient tensors and compares
the Wick factor $2!$ with a separate pairing of raw coefficients.
The saved output includes:
\begin{itemize}
\item 1,610 full seed-weight identities and 6,440 partial-weight checks
over five finite half-filled CAR systems and the two preparations;
\item 20,199 rational factorial estimates and 342 near-weight cases,
including asymmetric entries outside $\{0,1\}$;
\item 126 conditional Wick identities through degree six, with
repeated indices;
\item eight exact moment comparisons for rectangular tensors of
degrees one through three, checking all flattening orientations and
the coupled-chaos Jensen bound;
\item symbolic transfer expressions for the two preparation norms
and their two unnormalized $u^2$ responses.
\end{itemize}
These diagnostics do not replace the representation-theoretic and
probabilistic proofs above. They use no Haar Monte Carlo or floating
point extrapolation. The scientific JSON records parameters, software
versions and input SHA256 hashes. A separate read-only manifest binds
this draft, its PDF, source notes, build logs and parent evidence.
The successful driver is not rerun or overwritten during PDF production.

The use of tensor flattenings, Gaussian matrix inequalities and
decoupling is established methodology. Our project-specific step is
the control of every synthesis flattening by the near-Slater weight
estimate in the half-filled double-scaled regime. Equation
\eqref{conditional} is derived here to fix the repeated-index and
normalization conventions, rather than proposed as a new general
decoupling principle. Existing $q$-word rules supply the algebra in
\eqref{quad1}--\eqref{quad2}; the new input is the finite-$N$
transfer of exact pinching into these preparation states.

The next microscopic interface is the quadratic covariant generator
$J_2$ with selected $M$ in the same disorder, followed by the full
bilocal in this representation. The previously completed full-trace
bilocal four-point construction is retained. Interlaced pinched
moments, physical bulk spatial and mass generators and finite-$G$
dynamics remain open; no full de Sitter dictionary is claimed here.

\begin{thebibliography}{9}
\bibitem{chaos} A. S. Bandeira, K. Lucca, P. Nizi\'c-Nikolac and
R. van Handel, \emph{Matrix Chaos Inequalities and Chaos of
Combinatorial Type}, author manuscript dated 24 December 2024,
\S\S2.2--2.4 and Appendix A.2.
\url{https://web.math.princeton.edu/~rvan/chaosconf241224.pdf}.
\bibitem{tropp} J. A. Tropp,
\emph{Second-Order Matrix Concentration Inequalities},
\S\S6--7, especially the proof of Proposition 7.1.
\url{https://tropp.caltech.edu/papers/Tro16-Second-Order-Matrix-preprint.pdf}.
\bibitem{parent} A. Day and Chengheng,
\emph{Uniform Slater-orbit densities for linear selected preparations},
local project draft, 8 November 2026; and the preceding
\emph{Slater neutral-slope} and \emph{Slater covariant} stages.
Manuscripts, scientific outputs and provenance in this repository.
\end{thebibliography}
\end{document}
